\documentclass[11pt,a4paper]{article}
\usepackage[margin=2.25cm,headheight=15pt]{geometry}
\usepackage{fontspec}
\setmainfont{DejaVu Serif}
\setsansfont{DejaVu Sans}
\usepackage[spanish,es-noquoting]{babel}
\usepackage{amsmath,amssymb,booktabs,array,longtable,graphicx,xcolor,fancyhdr,enumitem,microtype}
\usepackage[colorlinks=true,linkcolor=teal,urlcolor=teal,citecolor=teal]{hyperref}
\definecolor{navy}{HTML}{0B1D32}\definecolor{teal}{HTML}{087E83}\definecolor{gold}{HTML}{AF642B}
\usepackage{titlesec}
\titleformat{\section}{\Large\bfseries\color{navy}}{\thesection}{.7em}{}
\titleformat{\subsection}{\large\bfseries\color{teal}}{\thesubsection}{.6em}{}
\pagestyle{fancy}\fancyhf{}\lhead{\footnotesize SUBATÓMICA · ADN HUMANO}\rhead{\footnotesize Proyecto de investigación}\cfoot{\thepage}
\newcommand{\I}{\mathrm I}\newcommand{\Te}{\mathrm{Te}}\newcommand{\Xe}{\mathrm{Xe}}
\newcommand{\pb}{\,\mathrm{pb}}\newcommand{\nm}{\,\mathrm{nm}}
\newcommand{\dato}[1]{\textcolor{teal}{\textsf{[DATO] #1}}}
\newcommand{\hip}[1]{\textcolor{gold}{\textsf{[MODELO] #1}}}
\setlength{\parskip}{5pt}\setlength{\parindent}{0pt}
\begin{document}
\begin{titlepage}
\color{navy}\vspace*{1.7cm}
{\sffamily\bfseries\Large UNIVERSIDAD NACIONAL DE COLOMBIA\par}
{\sffamily Departamento de Física · Física Subatómica\par}\vspace{2.2cm}
{\Huge\bfseries Cuando un núcleo decae dentro del ADN humano\par}\vspace{.55cm}
{\Large De la captura electrónica a las lesiones, la reparación y las mutaciones HPRT\par}
\vspace{1.4cm}\color{teal}\rule{\textwidth}{1.5pt}\color{navy}\vspace{.75cm}
{\large Guía científica y matemática para exposición tipo congreso\par}
{\large Artículo conductor: Whaley y Little (1990)\par}
\vfill
\textbf{Equipo:} Proyecto de Física Subatómica\par
\textbf{Fecha:} noviembre de 2026 · Bogotá, Colombia\par
\vspace{1cm}{\footnotesize Las figuras audiovisuales incluidas son esquemas y modelos propios; no son microfotografías ni datos de pacientes.}
\end{titlepage}
\begin{abstract}
¿Por qué el mismo radionúclido puede cambiar su efecto biológico al pasar de un medio extracelular al ADN? Resolvemos esta pregunta con la comparación experimental de Whaley y Little entre las líneas linfoblastoides humanas TK6 y SE30 expuestas a iododesoxiuridina marcada con $^{125}$I o $^{131}$I \cite{whaley}. Separamos cuatro niveles: el proceso nuclear que origina las radiaciones; la probabilidad de que el átomo decaiga mientras permanece en ADN; la geometría de las lesiones moleculares; y la selección experimental de mutantes HPRT. Derivamos las leyes de desintegración, un modelo de permanencia con riesgos competitivos, una regla geométrica de rotura doble y un modelo de destinos celulares. La única pendiente de mutación usada como cifra observada es $0.45\times10^{-6}$ mutantes inducidos por cGy para $^{131}$IUdR en TK6 \cite{whaley}. Las simulaciones de lesiones y reparación exploran hipótesis declaradas; no predicen el riesgo de una persona. La entrega incluye código, resultados y seis vídeos narrativos aptos para proyección.
\end{abstract}
\tableofcontents\clearpage

\section{El problema, su interés y una respuesta comprobable}
\subsection{Motivación}
La actividad radiactiva informa cuántas desintegraciones ocurren por segundo, pero no indica en qué parte de una célula se libera energía. Si un análogo de timidina, la iododesoxiuridina (IdUrd o IUdR), se incorpora al ADN durante la replicación, el átomo de yodo puede quedar próximo a la columna azúcar--fosfato. Si el átomo permanece fuera de la célula, la misma desintegración ocurre a otra distancia. La localización, la energía y la historia de reparación deben distinguirse. Whaley y Little compararon precisamente ambos escenarios en células humanas \cite{whaley}.

La aplicación inmediata es interpretar experimentos con radionúclidos dirigidos al núcleo y diseñar preguntas para radiofármacos. La comparación también enseña por qué una exposición ambiental a yoduro, una fuente de rayos X externa y un análogo de ADN marcado no se pueden declarar equivalentes por el simple hecho de contener radiación. El análisis conecta física nuclear, transporte de energía, biología molecular y lectura crítica de datos.

\subsection{Pregunta y objetivos}
\textbf{Pregunta:} ¿cómo afectan el modo de decaimiento y la distancia fuente--ADN la cadena entre desintegración y mutación HPRT en células humanas, y qué parte de esta cadena permite respaldar el experimento de Whaley y Little?

\textbf{Objetivos:} (1) derivar las leyes temporales de $^{125}$I y $^{131}$I; (2) distinguir captura electrónica, desexcitación atómica y beta menos; (3) estudiar por simulación cómo una hipótesis espacial modifica la probabilidad de una rotura doble; (4) separar lesión, mutación y pérdida celular; (5) calcular una extrapolación \emph{explícita} de la pendiente HPRT publicada; (6) comunicar resultados, supuestos y límites con vídeos y figuras reproducibles.

\subsection{Hipótesis y resultado principal}
\dato{El artículo observa mayores efectos de muerte y mutación cuando la IdUrd marcada se incorpora en ADN que cuando permanece en el medio, bajo sus condiciones de cultivo y sus controles TK6/SE30 \cite{whaley}.} \hip{Nuestra predicción cualitativa es que acercar la fuente a las dos hebras aumenta las lesiones agrupadas si la probabilidad de rotura decrece con la distancia.} Lo segundo es una consecuencia del modelo condicional, no una medición del artículo.

\section{Artículo guía: diseño, resultados y lectura crítica}
\textbf{Referencia íntegra:} Whaley, J. M. y Little, J. B., ``Efficient mutation induction by $^{125}$I and $^{131}$I decays in DNA of human cells'', \emph{Radiation Research} \textbf{123}, 68--74 (1990), DOI \href{https://doi.org/10.2307/3577659}{10.2307/3577659}; resumen abierto en \href{https://pubmed.ncbi.nlm.nih.gov/2371380/}{PubMed, PMID 2371380} \cite{whaley}. El artículo completo puede requerir acceso institucional; los resultados cuantitativos atribuidos aquí están explícitos en su resumen público.

TK6 es competente para la ruta de salvamento de timidina mediada por timidina cinasa (TK). SE30 tiene deficiencia en esa capacidad y permite comparar IdUrd que se incorpora con IdUrd cuya incorporación es deficiente. Se midieron muerte celular y mutantes en el locus \emph{hprt}. No se interpreta la comparación como un experimento de ADN purificado ni se atribuye a SE30 una incorporación estrictamente nula. El propio trabajo examinó explicaciones alternativas como sensibilidad basal distinta a radiación y sensibilización por IdUrd \cite{whaley}.

\begin{center}\small\begin{tabular}{p{.18\textwidth}p{.34\textwidth}p{.37\textwidth}}\toprule
Sistema & Fuente de radiación & Lectura \tabularnewline\midrule
TK6 + $^{125}$IIdUrd & Análogo incorporado en ADN & Efectos tóxicos y mutagénicos pronunciados, cualitativos en el resumen. \\
SE30 + $^{125}$IIdUrd & Incorporación deficiente; actividad en el medio & Prácticamente sin efecto bajo las condiciones descritas en el resumen. \\
TK6 + $^{131}$IIdUrd & Análogo incorporado en ADN & Pendiente HPRT inducida $0.45\times10^{-6}$ mutantes/cGy. \\
SE30 + $^{131}$IIdUrd & Actividad en el medio & Pocas mutaciones para dosis nucleares hasta $38$ cGy. \\
\bottomrule\end{tabular}\end{center}

\textbf{Precaución de lectura:} el valor $38$ cGy corresponde al límite de comparación indicado para células deficientes; aplicarlo a la pendiente de TK6 en la sección \ref{sec:hprt} es una \emph{extrapolación aritmética}, no un dato TK6 medido a esa dosis. No atribuimos al artículo valores numéricos de roturas dobles por cada desintegración de $^{125}$I ni una pendiente de mutación para $^{125}$I que no constan en el resumen.

Como apoyo geométrico, Martin y Haseltine estudiaron una deoxicitidina marcada con $^{125}$I en una posición definida de un fragmento de ADN: el daño radioquímico se concentró dentro de $15$--$20$ \AA{} ($1.5$--$2\nm$) y detectaron roturas hasta $70$ \AA{} ($7\nm$) \cite{martin}. Estas son distancias de daño en \emph{ese montaje experimental}, distintas del enunciado de Whaley y Little sobre depósito de gran parte de la energía de Auger a menos de $20\nm$ del sitio de desintegración \cite{whaley}. Ni un alcance de electrón ni la distribución de una rotura quedan completamente determinados por una sola longitud.

\section{Núcleo y curso: de la tasa a la actividad}
\subsection{Derivación completa de la ley exponencial}
Las diapositivas del curso \emph{Decays and scattering} introducen la probabilidad de decaimiento por unidad de tiempo y la densidad de estados. Escribimos $\lambda$ [día$^{-1}$] para la probabilidad instantánea de desintegración por átomo (el símbolo $\Gamma$ de ciertas diapositivas representa la misma tasa cuando su dimensión es tiempo$^{-1}$) \cite{curso}. La probabilidad de sobrevivir a un intervalo $dt$ es $1-\lambda dt+o(dt)$. Por ello, si $N(t)$ es el número medio de núcleos no desintegrados,
\begin{align}
N(t+dt)&=N(t)(1-\lambda dt)+o(dt),\\
N(t+dt)-N(t)&=-\lambda N(t)dt+o(dt),\\
\lim_{dt\to0}\frac{N(t+dt)-N(t)}{dt}&=-\lambda N(t).
\end{align}
Separamos variables, integramos desde $(0,N_0)$ hasta $(t,N(t))$, y usamos $\int dN/N=\ln N$:
\begin{align}
\frac{dN}{N}&=-\lambda\,dt, &
\int_{N_0}^{N(t)}\frac{dN'}{N'}&=-\lambda\int_0^t dt',\\
\ln\!\frac{N(t)}{N_0}&=-\lambda t, &N(t)&=N_0e^{-\lambda t}.
\end{align}
La semivida $T_{1/2}$ satisface $N(T_{1/2})=N_0/2$; entonces $1/2=e^{-\lambda T_{1/2}}$, $\ln(1/2)=-\lambda T_{1/2}$ y
\begin{equation}\boxed{\lambda=\frac{\ln2}{T_{1/2}},\qquad \tau=\lambda^{-1}=\frac{T_{1/2}}{\ln2}.}\end{equation}
La actividad es el \emph{número positivo de desintegraciones por tiempo}: $A(t)=-dN/dt=\lambda N(t)$, y $1$ Bq $=1$ s$^{-1}$. Si $\lambda$ se calcula por día, conviértase a s$^{-1}$ antes de informar actividad en Bq: $\lambda_{\mathrm{s}^{-1}}=\lambda_{\mathrm{día}^{-1}}/86400$. Un signo positivo delante de $dN/dt$ describiría una derivada negativa, no una actividad positiva.

\subsection{Ramas de decaimiento y amplitud cuántica}
Si existen rutas independientes con tasas $\lambda_k$, en un intervalo $dt$ la probabilidad total es $(\sum_k\lambda_k)dt$. De ahí $\lambda_{\mathrm{total}}=\sum_k\lambda_k$ y la razón de ramificación $b_k=\lambda_k/\lambda_{\mathrm{total}}$, con $\sum b_k=1$. Esto enlaza con la tasa de transición de Fermi $\lambda_{i\to f}\propto |\langle f|V|i\rangle|^2\rho(E_f)$ en las diapositivas \cite{curso}. El modelo temporal de supervivencia puede representarse por una amplitud $\psi(t)\propto e^{-iEt/\hbar}e^{-\lambda t/2}$; su módulo al cuadrado da $e^{-\lambda t}$. Aquí $\lambda$ es una tasa temporal; una anchura energética convencional es $\Gamma_E=\hbar\lambda$, que tiene unidades de energía. No igualamos una anchura en eV a una tasa en s$^{-1}$ sin este factor.

\section{$^{125}$I: captura electrónica y cascada Auger}
\subsection{Conservación de números cuánticos}
En la captura electrónica un protón nuclear capta un electrón ligado y se transforma en neutrón: $p+e^-\rightarrow n+\nu_e$ a escala de nucleones. El núcleo tiene $A=125$ antes y después, mientras $Z$ pasa de 53 a 52:
\begin{equation}{}^{125}_{53}\I+e^-_{\mathrm{ligado}}\longrightarrow{}^{125}_{52}\Te^*+\nu_e .\end{equation}
La carga se conserva en el proceso completo; el número leptónico del electrón inicial es $+1$ y el del neutrino saliente también. El estado nuclear excitado de $^{125}$Te se desexcita principalmente por conversión interna; sólo una fracción minoritaria emite directamente el fotón gamma de $35.49$ keV. IAEA informa $T_{1/2}=59.408$ d, captura electrónica $100\%$, gamma de $35.49$ keV con probabilidad $6.68\%$ y conversión interna aproximadamente $93.32\%$ desde ese estado \cite{iaea125}. Una conversión interna expulsa un electrón atómico usando energía de la transición nuclear: \emph{no} es el electrón previamente capturado.

\subsection{Balance de energía, sin simplificar electrones inexistentes}
Definimos la diferencia de masas \emph{atómicas neutras} en el estado fundamental como $Q_{EC}=[M_a(^{125}\I)-M_a(^{125}\Te)]c^2$. Sea $E_*$ la excitación nuclear de $^{125}$Te, $B_i$ la energía de enlace de la capa de la que se capturó el electrón y $T_R$ la energía de retroceso. A primer orden, despreciando pequeñas correcciones atómicas,
\begin{equation}E_{\nu,i}\simeq Q_{EC}-E_*-B_i-T_R.\label{eq:nu}\end{equation}
Explicación de cada sustracción: $E_*$ queda almacenada inicialmente en el núcleo hijo; $B_i$ representa la vacante atómica; $T_R$ es cinética del átomo hijo que garantiza momento conservado. El neutrino y el átomo tienen momentos opuestos a primer orden en un decaimiento de dos cuerpos, de modo que $T_R=p_R^2/(2M_R)$. El valor $35.49$ keV identifica la transición nuclear, \emph{no} es la energía de cada electrón Auger. Una vacante en una capa interna se rellena desde otra capa; la diferencia puede salir como rayo X característico o transferirse a un electrón Auger. Vacantes posteriores generan una cascada. Energías, multiplicidades y destinos requieren un espectro atómico; el código de esta entrega no inventa uno.

\subsection{¿Se queda el núcleo hijo unido al ADN?}
La iododesoxiuridina marcada incorpora \emph{el átomo padre} en una molécula del ADN antes del decaimiento. En el instante de captura, el núcleo hijo $^{125}$Te nace en el mismo lugar inicial, pero \emph{su posición inicial no garantiza que el nuevo elemento permanezca unido}. Cambian el número atómico, la configuración electrónica y los enlaces posibles; la relajación emite electrones y puede originar cambios químicos. Para estimar una escala de retroceso sólo por el neutrino, si su energía fuera $0.12$ MeV y se despreciara su masa, $p_\nu c\simeq0.12$ MeV. Con $M_Rc^2\simeq125(931.5)$ MeV,
\begin{equation}T_R\simeq\frac{(p_\nu c)^2}{2M_Rc^2}=\frac{(0.12\,\mathrm{MeV})^2}{2(125)(931.5\,\mathrm{MeV})}\simeq6.18\times10^{-8}\,\mathrm{MeV}\simeq0.062\,\mathrm{eV}.\end{equation}
Los $0.12$ MeV son un \emph{ejemplo de neutrino supuesto}; la energía real depende de la capa capturada y el estado final según la ecuación (\ref{eq:nu}). Aun si el retroceso nuclear en este ejemplo es pequeño, no permite concluir la estabilidad de la unión química tras la transmutación y la ionización múltiple. Para hablar de daño local basta que la fuente estuviera allí \emph{cuando} se produjo la cascada; nuestro modelo no simula el destino posterior del Te. En beta menos de $^{131}$I el producto es xenón: tampoco se conserva por definición el mismo enlace químico después de transformarse en otro elemento.

\subsection{Tiempo físico}
Usamos $T_{1/2}=59.4$ d redondeado para reproducir la animación. La tasa por día es $\lambda_{125}=\ln2/59.4=0.011669\,\mathrm{d}^{-1}$ aproximadamente. Si un millón de núcleos permanecen localizados durante 1 día sin pérdida biológica, el número \emph{medio} de desintegraciones es $10^6(1-e^{-0.011669})\simeq11\,601$. No todos los núcleos de una administración real quedan en el ADN ni siguen esta permanencia ideal.

\section{$^{131}$I: beta menos y un espectro continuo}
El mecanismo nuclear es $n\to p+e^-+\bar\nu_e$, que aumenta $Z$ de 53 a 54 y conserva $A=131$:
\begin{equation}{}^{131}_{53}\I\longrightarrow{}^{131}_{54}\Xe^*+e^-+\bar\nu_e .\end{equation}
Para una rama hacia una excitación hija $E_*$, definimos $Q_\beta=[M_a(^{131}\I)-M_a(^{131}\Xe)]c^2$ con masas atómicas neutras y $Q_b=Q_\beta-E_*$. Despreciando retroceso en un primer esquema,
\begin{align}Q_b&=T_e+E_{\bar\nu}+T_R,\\ E_{\bar\nu}&\simeq Q_b-T_e-T_R,\qquad 0\leq T_e\lesssim Q_b.\end{align}
El electrón no tiene energía fija porque el antineutrino participa en el reparto. Para una transición permitida ideal, la dependencia \emph{esquemática} es
\begin{equation}\frac{dN_b}{dT_e}\propto F(Z_f,T_e)\,p_e(T_e)\,\bigl(T_e+m_ec^2\bigr)\,(Q_b-T_e)^2\,C_b(T_e),\quad p_ec=\sqrt{T_e(T_e+2m_ec^2)}.\end{equation}
Aquí $F$ corrige la interacción coulombiana, $C_b$ incorpora forma y reglas de transición, y $b$ es una rama nuclear concreta. El espectro completo sería la suma ponderada $dN/dT_e=\sum_b b_b(dN_b/dT_e)$; no se declara que la fórmula simplificada dé el espectro real completo de $^{131}$I. Su semivida es $8.021$ d según ICRP \cite{icrp}; la emisión gamma principal próxima a $364.5$ keV y la rama beta con extremo próximo a $606$ keV se documentan para este isótopo en IAEA \cite{iaea131}. $606$ keV es energía \emph{máxima de una rama}, no energía única ni energía media.

La existencia de un electrón de mayor recorrido no implica igual probabilidad de daño por desintegración incorporada: localización, cascadas secundarias y geometría importan. La diapositiva de detección del curso discute pérdida de energía y alcance; sus ejemplos de pico de Bragg para hadrones no se transfieren sin más a los electrones beta, sujetos a dispersión y múltiples interacciones \cite{deteccion}.

\section{Un átomo puede abandonar el ADN antes de decaer}
Sea $n(t)$ la población media de átomos aún en ADN. Además del decaimiento con tasa $\lambda$, supongamos \emph{como hipótesis} una salida exponencial con tasa biológica $k$ [día$^{-1}$]. En $dt$, ambas vías son eventos competitivos: $dn=-(\lambda+k)n\,dt$, por lo que
\begin{equation}n(t)=n_0e^{-(\lambda+k)t}.\end{equation}
Las desintegraciones \emph{dentro del ADN} durante $dt$ son $dD_{ADN}=\lambda n(t)dt$. Integrando desde 0 hasta $t$:
\begin{align}
D_{ADN}(t)&=\int_0^t\lambda n_0e^{-(\lambda+k)s}\,ds\\
&=\lambda n_0\left[\frac{-e^{-(\lambda+k)s}}{\lambda+k}\right]_0^t\\
&=n_0\frac{\lambda}{\lambda+k}\left(1-e^{-(\lambda+k)t}\right).
\end{align}
Dividiendo por $n_0$ se obtiene $P(D_{ADN}\text{ antes de }t)$. Si $t\to\infty$, el término exponencial tiende a cero, y $P=\lambda/(\lambda+k)$. Sea $\tau_b$ el tiempo medio de residencia: $k=1/\tau_b$, \emph{no} $\ln2/\tau_b$. Con la hipótesis $\tau_b=2$ d, $k=0.5$ d$^{-1}$ y las tasas anteriores, obtenemos $P_{125}=0.0228$ y $P_{131}=0.1474$ aproximadamente. Estos porcentajes \emph{no} describen la biocinética de las células del artículo: ilustran el peso de una tasa competidora. \textbf{Chequeo:} $k=0$ da $P=1$; $\lambda\to0$ con $k>0$ da $P\to0$.

\section{Geometría del ADN y modelo explícito de daño}
\subsection{Escala y coordenadas}
Para una sección ideal de $n=60$ pares de bases usamos radio de columna $R=1.0$ nm y avance axial $h=0.34$ nm por par (vuelta aproximada de 10 pb = 3.4 nm). Estas cifras fijan una \emph{hélice ilustrativa}, no la estructura atomística de un locus HPRT. Definimos $j=0,\ldots,n-1$, $j_c=\lfloor n/2\rfloor$ y $\theta_j=2\pi(j-j_c)/10$. Las dos columnas son
\begin{equation}\mathbf x_{A,j}=(R\cos\theta_j,R\sin\theta_j,h(j-j_c)),\qquad
\mathbf x_{B,j}=(-R\cos\theta_j,-R\sin\theta_j,h(j-j_c)).\end{equation}
La fuente se ubica en $\mathbf s=(R+d,0,0)$: $d=0$ corresponde al borde de la hebra A en el par central, y $d>0$ mueve la fuente hacia afuera. La distancia euclidiana de fuente a cada sitio es $r_{X,j}=\|\mathbf x_{X,j}-\mathbf s\|$, para $X\in\{A,B\}$. \emph{No} se sustituyen pares de bases por nanómetros sin la conversión $h=0.34\nm/\mathrm{pb}$, y para pares opuestos a igual $j$ queda además separación transversal.

\subsection{De un peso espacial a una Bernoulli de rotura}
Introducimos el peso \emph{fenomenológico} $q_{X,j}=\exp(-r_{X,j}/\ell)$, donde $\ell=1.5$ nm es un parámetro \emph{elegido para explorar}, no ajustado al alcance observado en Martin y Haseltine. Supongamos oportunidades de lesión locales Poisson de media $\mu q_{X,j}$, con $\mu=0.48$ como escenario. La probabilidad de \emph{cero} oportunidades es $\exp(-\mu q_{X,j})$. Por complemento, la probabilidad de al menos una rotura del sitio es
\begin{equation}\boxed{p_{X,j}=1-\exp[-\mu\exp(-r_{X,j}/\ell)].}\label{eq:ssb}\end{equation}
Cada evento simulado genera una variable uniforme $U_{X,j}\sim\mathcal U(0,1)$ y marca rotura si $U_{X,j}<p_{X,j}$. Las Bernoulli se tratan como independientes \emph{condicionadas a} $d,\mu,\ell$. Una cascada real tiene correlaciones entre electrones y química acuosa; la independencia es un límite del modelo. No se asigna a cada oportunidad una energía en eV ni se convierte la expresión en un espectro Auger.

\subsection{Criterio DSB y estimación Monte Carlo}
La regla operativa usada aquí llama ``DSB'' a la existencia de dos roturas, una en A y otra en B, con $|j_A-j_B|\leq10$ pb. La ventana de 10 pb es una \emph{convención algorítmica declarada}; variarla cambia la clasificación. Para $N_{MC}$ eventos independientes con indicador $I_m(d)\in\{0,1\}$:
\begin{equation}\widehat p_{DSB}(d)=\frac{1}{N_{MC}}\sum_{m=1}^{N_{MC}}I_m(d),\qquad
\widehat{\mathrm{SEM}}=\sqrt{\frac{\widehat p(1-\widehat p)}{N_{MC}}}.\end{equation}
El error estándar binomial estima \emph{variación de muestreo con parámetros fijos}. No incluye incertidumbre sobre $\ell$, $\mu$, la geometría o los mecanismos de reparación. Se emplea una semilla fija (2026) y números aleatorios comunes por escenario $d$, lo que facilita reproducir diferencias. Cada evento está condicionado a una desintegración en una posición elegida: $\widehat p_{DSB}$ \emph{no} es frecuencia medida de DSB por decaimiento del artículo.

\subsection{Control analítico: rotura en cualquier sitio}
Antes de estudiar las parejas DSB, existe una prueba cerrada del simulador. Denotemos por $B_{X,j}=0$ que el sitio no se rompió. Como las variables Bernoulli son independientes \emph{dentro de esta hipótesis}, la probabilidad de ninguna rotura en ningún sitio es el producto de todas las probabilidades complementarias:
\begin{align}
P(\text{sin roturas})&=\prod_{X\in\{A,B\}}\prod_{j=0}^{59}(1-p_{X,j})\\
&=\prod_{X,j}\exp[-\mu\exp(-r_{X,j}/\ell)]\\
&=\exp\left[-\mu\sum_{X,j}\exp(-r_{X,j}/\ell)\right].
\end{align}
Pasamos de producto a exponencial de suma porque $\prod e^{a_j}=e^{\sum a_j}$; tomando el complemento se obtiene $P(\text{alguna rotura})=1-\exp[-\mu\sum_{X,j}q_{X,j}]$. Con la fuente en $d=0$, la fórmula da $0.99693$ y el código obtuvo $0.99713$ en $15\,000$ eventos: la diferencia $0.00020$ es pequeña frente a su error binomial aproximado. Esto verifica el muestreo de roturas individuales \emph{según el modelo}, no valida la biofísica de los electrones reales.

Una pareja opuesta concreta $(i,j)$, con $|i-j|\leq10$, tiene probabilidad $p_{A,i}p_{B,j}$ bajo la misma independencia. Al sumar esas probabilidades se obtiene el \emph{número medio de parejas}, no la probabilidad de al menos una DSB: un evento puede contener varias parejas. La desigualdad de la unión sólo asegura $P(DSB)\leq\sum_{|i-j|\leq10}p_{A,i}p_{B,j}$. Esta distinción explica por qué el algoritmo usa la función indicadora ``existe alguna pareja'', en vez de sumar cuentas de parejas como si fueran eventos separados.

\subsection{Resultados del escenario base y sensibilidad}
Los resultados reproducibles de \texttt{resultados/resumen\_modelo.json} para $N_{MC}=15\,000$, $\mu=0.48$, $\ell=1.5$ nm, semilla 2026, son:
\begin{center}\small\begin{tabular}{rrrr}\toprule
$d$ [nm]& $P(\text{cualquier rotura})$ & $\widehat P(DSB)$ & SEM$_{DSB}$ \\ \midrule
0&0.99713&0.87000&0.00275\\
1&0.98167&0.68847&0.00378\\
2&0.90880&0.41560&0.00402\\
4&0.55587&0.07180&0.00211\\
7&0.12580&0.00200&0.00036\\ \bottomrule
\end{tabular}\end{center}
El error de muestreo al 95\% se aproximaría por $\widehat p\pm1.96\,SEM$ cuando la aproximación normal es adecuada; para probabilidades cercanas a cero conviene un intervalo binomial exacto o de Wilson. Ninguno de estos intervalos cubriría automáticamente la incertidumbre de $\ell$ y $\mu$.

Hicimos un barrido intencional de \emph{hipótesis}, con $6\,000$ eventos por punto y la misma semilla. A $d=0$, si $\mu$ baja a $0.24$, $\widehat P(DSB)=0.531$; si sube a $0.96$, pasa a $0.993$. Con $\mu=0.48$ y $\ell=1.0$ nm, resulta $0.593$; con $\ell=2.0$ nm, $0.963$. En $d=4$ nm esas mismas alternativas producen $0.0217$, $0.2147$, $0.0032$ y $0.3003$, respectivamente. La variación excede por mucho los errores de muestreo: el modelo no fija una probabilidad biológica universal. El patrón \emph{dentro de cada escenario} es que la fuente más cercana da más parejas de roturas; la magnitud debe establecerse experimentalmente o con transporte microscópico validado.

\begin{figure}[htbp]\centering\includegraphics[width=.92\textwidth]{../animaciones/04_roturas_agrupadas_portada.png}\caption{La regla sobre dos hebras y un barrido espacial con parámetros del modelo. La gráfica es un resultado simulado; no son datos experimentales.}\end{figure}

\subsection{De energía a dosis: qué representa un gray}
La energía absorbida $E_{abs}$ dividida por una masa $m$ define $D=E_{abs}/m$ y $1\,\mathrm{Gy}=1\,\mathrm{J\,kg^{-1}}$. Para un volumen de agua ideal esférico, $m=\rho(4\pi/3)r^3$. Si \emph{hipotéticamente} se depositara $1$ keV entero en una esfera de agua de radio $2$ nm, se obtendría $E=10^3(1.602176634\times10^{-19})=1.602176634\times10^{-16}$ J, $m=(10^3\,\mathrm{kg\,m^{-3}})(4\pi/3)(2\times10^{-9}\,\mathrm m)^3=3.3510\times10^{-23}$ kg y $D\simeq4.78\times10^6$ Gy para \emph{esa masa y ese único evento supuesto}. Ese valor no es una dosis terapéutica, no significa que un electrón real deposite íntegramente 1 keV ahí, y no se compara con una dosis celular promedio sin especificar el volumen y el reparto de energía. Precisamente muestra por qué la escala del volumen debe escribirse cuando se calcula un gray.

\section{Después del daño: tres destinos que compiten}
Sea $L(t)$ la fracción de lesiones aún presentes y $R(t)$, $M(t)$, $C(t)$ las fracciones acumuladas que ya han pasado a reparación, mutación fijada y pérdida celular. El tiempo está en horas y las tasas \emph{hipotéticas} $k_r=0.20$, $k_m=0.02$, $k_c=0.03$ h$^{-1}$. Este $k_m$ es \emph{otro parámetro} distinto de la pendiente de mutación observada en HPRT. Con $K=k_r+k_m+k_c=0.25$ h$^{-1}$:
\begin{align}
\frac{dL}{dt}&=-KL,&L(0)&=1,\\
\frac{dR}{dt}&=k_rL,&\frac{dM}{dt}&=k_mL,&\frac{dC}{dt}&=k_cL.
\end{align}
Integramos: $L(t)=e^{-Kt}$ y $R(t)=k_r\int_0^t e^{-Ks}ds=(k_r/K)(1-e^{-Kt})$; sustituyendo $k_r$ por $k_m,k_c$ obtenemos $M(t),C(t)$. Sumando,
\begin{equation}L+R+M+C=e^{-Kt}+\frac{k_r+k_m+k_c}{K}(1-e^{-Kt})=1.\end{equation}
En el límite $t\to\infty$ el destino simulado se reparte $80\%$, $8\%$, $12\%$. Son fracciones \emph{asumidas por las tasas}, no prevalencias reales de lesiones ni una calibración humana. ``Mutación fijada'' exige supervivencia y cambios heredables en descendientes celulares: la animación representa la distinción, no reemplaza experimentos de linaje y secuenciación.

\section{Del locus HPRT a un conteo interpretable}\label{sec:hprt}
Una rotura observada en ADN no demuestra una mutación heredable; una mutación HPRT se evalúa con una selección y la frecuencia de mutantes entre células clonables. El experimento del artículo registra mutantes en el locus \emph{hprt} y distingue la muerte celular de esta variable \cite{whaley}.
\textbf{Dos enzimas, dos funciones:} la timidina cinasa (TK) controla la incorporación del análogo IdUrd; la hipoxantina--guanina fosforribosiltransferasa (HPRT) participa en la vía de rescate de purinas y permite seleccionar células resistentes a 6-tioguanina cuando se pierde la función correspondiente. TK6 y SE30 comparan la primera ruta; la medición de mutantes HPRT explora otro locus. No sería correcto afirmar que ``SE30 tiene mutado HPRT'' por ser deficiente en TK. HPRT1 humano se encuentra en el cromosoma X \cite{gene}; O'Neill, Sullivan y Albertini describieron un ensayo de clonación con 6-tioguanina y corrección por capacidad de clonación en células humanas \cite{oneill}. Este apoyo metodológico no demuestra que el procedimiento exacto de cultivo de Whaley y Little usara idénticas concentraciones o tiempos.

Si $n$ células se siembran por pocillo y cada célula tiene probabilidad efectiva $f\,\varepsilon$ de producir un clon resistente, donde $f$ es la frecuencia de mutantes y $\varepsilon$ es su eficiencia de clonación bajo selección, entonces para eventos raros la cantidad de clones resistentes por pocillo se modela como Poisson con media $nf\varepsilon$. La probabilidad de un pocillo sin clon es $P_0=e^{-nf\varepsilon}$. Aplicando logaritmo a ambos lados: $\ln P_0=-nf\varepsilon$, de donde
\begin{equation}\boxed{f=-\frac{\ln P_0}{n\varepsilon}.}\end{equation}
Esta expresión ilustra por qué contar sólo pocillos positivos no basta: deben conocerse el número sembrado y la eficiencia de clonación, y debe permitirse expresión fenotípica. El programa de esta entrega parte \emph{ya} de un número de células clonables y usa Poisson para ilustrar repeticiones; no reconstruye los datos crudos del artículo.

Para modelar un conteo simplificado con $n$ células clonables y frecuencia de fondo $f_0$, usamos
\begin{equation}f_{TK6}(D)=f_0+\alpha D,\qquad \alpha=0.45\times10^{-6}\,\mathrm{cGy}^{-1},\qquad D\ \text{en cGy}.\end{equation}
La pendiente $\alpha$ corresponde a $^{131}$IIdUrd en TK6 del artículo; el fondo $f_0$ no está informado aquí y en la simulación se \emph{elige} $10\times10^{-6}$. El exceso de frecuencia a $D=38$ cGy es
\begin{align}
\Delta f&=\alpha D=(0.45\times10^{-6}\,\mathrm{cGy}^{-1})(38\,\mathrm{cGy})\\
&=17.1\times10^{-6}=17.1\text{ mutantes inducidos por millón de células clonables},\\
38\,\mathrm{cGy}&=0.38\,\mathrm{Gy},\qquad \alpha=45\times10^{-6}\,\mathrm{Gy}^{-1}.
\end{align}
Se cambió la unidad porque $1$ Gy $=100$ cGy; para mantener $\alpha D$ idéntico, el valor numérico de $\alpha$ por Gy se multiplica por $100$. La cifra $17.1$ es una \textbf{extrapolación de TK6} en un nivel de dosis mencionado para SE30, no una medición directa de TK6 a 38 cGy y tampoco una tasa de mutación del genoma completo. El modelo opcional de réplicas toma $Y_0\sim\mathrm{Poisson}(nf_0)$ y $Y_D\sim\mathrm{Poisson}(n(f_0+\alpha D))$; por independencia, $E[Y_D-Y_0]=n\alpha D$ y $\mathrm{Var}(Y_D-Y_0)=n(2f_0+\alpha D)$. El número de células \emph{clonables} se presupone, no se infiere de supervivencia.

\section{Resultados reproducibles y lectura de las seis animaciones}
Ejecutar \texttt{python codigo/modelo.py} produce \texttt{resultados/resumen\_modelo.json}. Ejecutar \texttt{python codigo/animar.py} produce seis MP4 de 1280$\times$720 a 20 fps, un vídeo unido y seis portadas PNG. Cada vídeo dura aproximadamente 7 s; se recomienda pausar el vídeo y explicar la variable indicada. Los cuadros ilustran mecanismos, no una filmación real ni trayectorias calculadas con Geant4.

\begin{longtable}{p{.24\textwidth}p{.32\textwidth}p{.35\textwidth}}\toprule Escena & Lo que cambia realmente & Qué explicar frente al público \\ \midrule\endhead
01 Captura y Auger & Representación de I-125 a Te-125 con vacante y cascada. & Conservación de $A$ y $Z$, captura de electrón ligado, excitación y Auger; el movimiento audiovisual no es un tiempo físico. \\
02 Semividas & $N/N_0=e^{-\lambda t}$; ejemplo de salida con $k=0.5$ d$^{-1}$. & Dos escalas temporales, fracción que decae antes de salir, diferencia entre $T_{1/2}$ y residencia. \\
03 Distancia & $q(r)=e^{-r/\ell}$, con $\ell=1.5$ nm hipotético. & El perfil sirve al modelo; no es el alcance medido de un electrón ni una curva experimental. \\
04 Dos hebras & Simulación de Eq. (\ref{eq:ssb}) y ventana de 10 pb. & Rotura simple frente a rotura doble; frecuencia condicional y error de muestreo. \\
05 Destinos & $L,R,M,C$ con suma unitaria. & Daño, reparación, fijación y pérdida son variables distintas; tasas ilustrativas. \\
06 HPRT humano & TK6/SE30, pendiente publicada de I-131. & ¿Qué se midió? ¿Qué se extrapoló? ¿Qué no se puede inferir sobre una persona? \\
\bottomrule\end{longtable}

\section{Vínculo preciso con el programa}
Las diapositivas proporcionadas de \emph{Decays and scattering} tratan estados inestables, tasa, semivida, actividad, modos y razón de ramificación: aparecen en las secciones 3--5 \cite{curso}. La presentación \emph{Particle Detection--part 2} introduce pérdida de energía y alcance de partículas en materia: motiva nuestra discusión sobre deposición de energía y límites de un modelo de electrones \cite{deteccion}. La física de interacciones débiles aparece en la transformación nuclear beta/captura, mientras la relajación electromagnética permite distinguir rayos X, gamma y Auger. No se usan QCD, modelo de capas, Proca o diagramas de Feynman como adornos: no son necesarios para responder esta pregunta. El ensayo HPRT entra como medición biológica posterior y no sustituye una medida de energía de detector.

\section{Utilidad, analogías prudentes y limitaciones}
\textbf{Utilidad científica:} seleccionar blancos subcelulares para estudiar emisores Auger; priorizar medidas de incorporación, distribución y permanencia antes de convertir Bq en efecto; distinguir un biomarcador de daño (roturas) de un desenlace tardío (mutación). \textbf{Símil de actualidad:} una radioterapia con radiofármaco dirigido al núcleo exige saber dónde queda cada fuente; las exposiciones por accidentes nucleares requieren otras vías de entrada, metabolismo y distribución. Los resultados de IdUrd en cultivo no cuantifican las consecuencias sanitarias de Chernóbil, bombas atómicas o radioyodo convencional, ni se extrapolan a personas sin datos de órganos y dosis.

\textbf{Límites del estudio:} (i) las líneas celulares humanas no son una cohorte de pacientes; (ii) la actividad y la dosis al núcleo no identifican por sí solas una secuencia genómica mutada; (iii) la hélice representada contiene 60 pb ideales, no el genoma humano; (iv) el modelo espacial usa un núcleo de probabilidad elegido, independencia de sitios y un criterio operacional DSB; (v) la tasa de residencia, las tasas celulares y el fondo HPRT del código son escenarios; (vi) los vídeos muestran un relato de mecanismos y no mediciones en tiempo real. El experimento de LeMotte y Little encontró aproximadamente una rotura doble por desintegración $^{125}$I bajo condiciones congeladas \cite{lemotte}; \emph{esa} medida no valida las probabilidades numéricas de nuestro código para una hélice aislada. La corrección de Schmitz y colaboradores subrayó además que una rotura doble no se convierte automáticamente en aberración cromosómica: sus datos revisados dieron en promedio unas 250 desintegraciones por aberración bajo sus condiciones experimentales \cite{corrigendum}.

\section{Conclusión defendible}
Las ecuaciones nucleares y la ley de decaimiento describen \emph{cuándo y qué} se emite. La geometría y la permanencia modifican \emph{dónde} se deposita energía. Reparación, supervivencia y selección experimental determinan \emph{qué efecto final se observa}. En cultivos humanos TK6/SE30, Whaley y Little encontraron una diferencia clara entre radionúclidos incorporados y actividad extracelular \cite{whaley}. Nuestro modelo muestra una vía consistente para pensar esa diferencia y expone los datos necesarios para verificarla con mayor fidelidad, sin reclamar una predicción clínica o genómica que los datos no sustentan.

\appendix\section{Guía de defensa breve: preguntas y respuestas}
\begin{enumerate}[leftmargin=*]
\item \textbf{¿Por qué captura electrónica no es beta menos?} En EC, un protón capta un electrón ligado y $Z$ baja; en beta menos, un neutrón produce electrón y antineutrino y $Z$ sube.
\item \textbf{¿$^{125}$I emite electrones monoenergéticos?} La captura deja vacantes y puede iniciar conversiones y cascadas de energía diversa; un dibujo con seis rayos no representa seis electrones con la misma energía.
\item \textbf{¿Por qué 38 cGy dan 17.1 mutantes?} Multiplicamos la pendiente TK6 del artículo por 38 cGy como \emph{ejercicio extrapolado}; 38 cGy se menciona para el control SE30.
\item \textbf{¿Si un DSB ocurre, ya hay mutación?} No. Debe intervenir reparación, supervivencia y fijación; un daño puede desaparecer o matar la célula.
\item \textbf{¿Qué validarían más experimentos?} Dependencia de daño con posición del radionúclido, espectro electrónico real, dosimetría celular, supervivencia y secuenciación de clones supervivientes.
\end{enumerate}

\begin{thebibliography}{99}
\bibitem{whaley} J. M. Whaley y J. B. Little, ``Efficient mutation induction by $^{125}$I and $^{131}$I decays in DNA of human cells'', \emph{Radiat. Res.} 123 (1990) 68--74. \href{https://pubmed.ncbi.nlm.nih.gov/2371380/}{PMID 2371380}. DOI: \href{https://doi.org/10.2307/3577659}{10.2307/3577659}.
\bibitem{martin} R. F. Martin y W. A. Haseltine, ``Range of radiochemical damage to DNA with decay of iodine-125'', \emph{Science} 213 (1981) 896--898. \href{https://pubmed.ncbi.nlm.nih.gov/7256283/}{PMID 7256283}. DOI: \href{https://doi.org/10.1126/science.7256283}{10.1126/science.7256283}.
\bibitem{lemotte} P. K. LeMotte y J. B. Little, ``DNA damage induced in human diploid cells by decay of incorporated radionuclides'', \emph{Cancer Res.} 44 (1984) 1337--1342. \href{https://pubmed.ncbi.nlm.nih.gov/6704953/}{PMID 6704953}.
\bibitem{oneill} J. P. O'Neill, L. M. Sullivan y R. J. Albertini, ``In vitro induction, expression and selection of thioguanine-resistant mutants with human T-lymphocytes'', \emph{Mutat. Res.} 240 (1990) 135--142. \href{https://pubmed.ncbi.nlm.nih.gov/2300074/}{PMID 2300074}. DOI: \href{https://doi.org/10.1016/0165-1218(90)90017-V}{10.1016/0165-1218(90)90017-V}.
\bibitem{gene} NCBI Gene, \emph{HPRT1 hypoxanthine phosphoribosyltransferase 1 [Homo sapiens]}, Gene ID 3251, \url{https://www.ncbi.nlm.nih.gov/gene/3251}.
\bibitem{iaea125} IAEA, \emph{IAEA-TECDOC-1340}, ficha de $^{125}$I, sección 75. \url{https://www-pub.iaea.org/MTCD/publications/PDF/te_1340_web.pdf}.
\bibitem{icrp} ICRP, \emph{Publication 137}, Occupational Intakes of Radionuclides, parte 3 (2017), tabla de semividas de yodo. \url{https://doi.org/10.1177/0146645317734963}.
\bibitem{iaea131} IAEA, \emph{IAEA-TECDOC-1608}, sección 12.3, propiedades de $^{131}$I. \url{https://www-pub.iaea.org/MTCD/Publications/PDF/te_1608_web.pdf}.
\bibitem{corrigendum} S. Schmitz et al., ``Corrigendum to `Chromosome aberrations induced by the Auger electron emitter $^{125}$I''', \emph{Mutat. Res.} 800--801 (2016) 46. \href{https://doi.org/10.1016/j.mrgentox.2016.03.003}{DOI 10.1016/j.mrgentox.2016.03.003}.
\bibitem{curso} C. Sandoval, \emph{Introducción a la física subatómica: Decays and scattering}, diapositivas de curso, Universidad Nacional de Colombia, 2025-1 (material aportado por el estudiante).
\bibitem{deteccion} C. Sandoval, \emph{Introducción a la física subatómica: Particle Detection--part 2}, diapositivas de curso, Universidad Nacional de Colombia, 2025-1 (material aportado por el estudiante).
\end{thebibliography}
\end{document}
